\documentclass[12pt,a4paper]{article}
\usepackage{amsmath,booktabs,graphicx,hyperref,geometry,setspace,authblk}
\geometry{margin=2.5cm}
\onehalfspacing

\title{A cross-cohort workflow for detecting microbiome biomarker
non-portability under compositional cohort structure}

\author[1]{[Author names]}
\affil[1]{[Institution]}

\date{\today}

\begin{document}
\maketitle

% ============================================================
\begin{abstract}
\textbf{Background:}
Microbiome biomarkers are often evaluated using pooled case--control
comparisons or random cross-validation. These approaches can overestimate
biomarker performance when disease labels are embedded within
cohort-specific compositional structure.

\textbf{Methods:}
We developed and validated a cross-cohort analytical workflow to detect
microbiome biomarker non-portability under compositional cohort-structured
heterogeneity. The workflow integrates centred log-ratio (CLR)
transformation, Aitchison-space variance partitioning, a cohort--disease
dominance ratio (CDDR), disease-effect vector geometry, leave-one-cohort-out
(LOCO) machine-learning validation, biomarker portability gap (BPG)
estimation, domain-correction feasibility assessment, and feature-stability
analysis. We applied it to three public Parkinson's disease gut microbiome
cohorts from Finland, Malaysia, and the United States, comprising 683
individuals and 62 harmonised genus-level CLR features. Genus-level
p-values were corrected for multiple comparisons using the
Benjamini--Hochberg (BH) false-discovery-rate procedure applied separately
within each hypothesis family.

\textbf{Results:}
Cohort structure dominated the compositional space. PC1 explained 61.2\%
of total variance; cohort identity accounted for 75.05\% of
Aitchison-space variance ($R^2 = 0.7505$) while disease status accounted
for 0.29\% ($R^2 = 0.0029$), yielding a CDDR of 258.7. The disease-effect
vector was nearly orthogonal to PC1 (cosine similarity = 0.067;
angle = 86.1$^\circ$). In genus-level linear models, 20 of 62 genera showed
a cohort-adjusted PD main effect surviving BH correction ($q < 0.05$), and
25 of 62 genera showed significant PD-by-cohort interaction ($q < 0.05$).
Thirty-one genera exhibited significant cross-cohort heterogeneity in the PD
effect (Cochran $Q$; $q < 0.05$), and only 12 genera showed a consistent
pooled direction of effect across cohorts ($q_{\mathrm{FE}} < 0.05$). Pooled
cross-validation substantially overestimated biomarker performance: AUROC
values reached 0.810, 0.881, and 0.919 for elastic-net logistic regression,
random forest, and XGBoost. Under LOCO validation, mean AUROC decreased to
0.600, 0.515, and 0.547, respectively, producing AUROC portability gaps of
0.210--0.371. Feature importance was unstable across held-out-cohort training
regimes (mean top-10 overlap: 2.0--3.0 features). Transductive
cohort-centring produced degenerate classifiers (AUROC = 0.500 for all
three models), consistent with the interpretation that cohort-structured
variation drove pooled predictive performance.

\textbf{Conclusions:}
Apparent microbiome biomarker performance can arise from cohort-structured
compositional variation rather than portable disease signal. The proposed
workflow identifies this failure mode by combining compositional geometry,
cohort-held-out prediction, portability-gap estimation, and feature-stability
analysis, and provides a practical validation strategy for microbiome
biomarker studies before candidate signatures are interpreted as
transportable across cohorts.

\noindent\textbf{Keywords:}
microbiome biomarkers; biomarker portability; cohort effects;
compositional data analysis; cross-cohort validation; machine learning;
Aitchison geometry; Parkinson's disease
\end{abstract}

% ============================================================
\section{Introduction}
\label{sec:intro}

Microbiome studies increasingly use high-throughput sequencing and
machine-learning models to identify disease-associated microbial signatures.
A recurring limitation is that signatures discovered in one dataset often
show reduced reproducibility in independent cohorts. This problem is usually
discussed as a matter of technical variation, sample size, sequencing
protocol, or clinical heterogeneity. However, in microbiome data the problem
is also geometric: relative abundance profiles are compositional, and
cohort-specific community structure can dominate the feature space in which
disease effects are estimated.

This creates a specific failure mode for microbiome biomarker discovery. A
classifier may perform well under pooled random cross-validation because
training and test folds share cohort-specific compositional structure. The
same classifier may fail under cohort-held-out validation, where the disease
signal must transport to a different microbiome background. In this setting,
statistical association does not imply portability. A genus-level disease
signature can be detectable within cohorts yet unstable across cohorts if the
disease-effect vector is small, cohort-conditioned, or misaligned with the
dominant compositional gradient.

We therefore developed a cross-cohort workflow to diagnose microbiome
biomarker non-portability. The workflow combines centred log-ratio
transformation, Aitchison-space variance partitioning, a cohort--disease
dominance ratio, disease-effect vector geometry, leave-one-cohort-out
prediction, biomarker portability gap estimation, domain-correction
feasibility assessment, and feature-stability analysis. Parkinson's disease
(PD) gut microbiome cohorts provide a relevant test case because PD has
repeatedly been associated with gut microbiome changes, while reported
microbial taxa vary considerably across
studies~\cite{Sampson2016,Bedarf2017,Barichella2019}. The aim was not to
nominate a universal PD microbiome biomarker, but to evaluate whether
apparent disease-associated microbial signatures remain portable across
cohort backgrounds.

% ============================================================
\section{Materials and Methods}
\label{sec:methods}

\subsection{Datasets and harmonisation}

Three publicly available 16S rRNA amplicon sequencing cohorts were used.
The Finland cohort comprised 148 individuals (74 PD, 74 controls). The
Malaysia cohort comprised 198 individuals (101 PD, 97 controls). The United
States cohort comprised 337 individuals (204 PD, 133 controls). Total: 683
individuals. Raw sequence data were processed through a unified MOTHUR
pipeline to canonical genus-level taxonomy. Genera were retained if present
in all three cohorts at the same taxonomic annotation; 62 genus-level
features were included in the harmonised feature matrix. Metadata variables
were PD status (binary) and cohort membership; both were complete with no
missing values.

\subsection{Centred log-ratio transformation}

Compositional genus count matrices were CLR-transformed after addition of a
pseudocount of 1 to handle zeros~\cite{Aitchison1982}. For each sample,
$x_{j,\mathrm{CLR}} = \log(x_j + 1) - \overline{\log(\mathbf{x}+1)}$,
where the mean is taken over all $p = 62$ genera. The CLR transformation
was applied globally to the pooled 683-sample matrix.

\subsection{Aitchison-space variance partitioning}

PERMANOVA (permutational multivariate analysis of variance) was applied to
the Aitchison distance matrix (Euclidean distance on CLR-transformed data)
using \texttt{adonis2} from the \texttt{vegan} package (999
permutations)~\cite{Anderson2001,Oksanen2022}. Both a sequential model
($\sim\mathrm{Cohort} + \mathrm{PD}$) and a marginal (by-margin) model were
fitted. The cohort--disease dominance ratio (CDDR) was defined as
$\mathrm{CDDR} = R^2_{\mathrm{Cohort}} / R^2_{\mathrm{PD}}$ from the
marginal model.

\subsection{Disease-effect vector geometry}

PCA was performed on the CLR feature matrix. The disease-effect vector
$\mathbf{d}$ was computed as the mean CLR profile of PD subjects minus the
mean CLR profile of control subjects (both pooled across cohorts). Alignment
of $\mathbf{d}$ with PC1 was measured as $\cos\theta = (\mathbf{d} \cdot
\mathbf{v}_1) / (\|\mathbf{d}\|\,\|\mathbf{v}_1\|)$, where $\mathbf{v}_1$ is
the PC1 loading vector. The angle $\theta = \arccos(|\cos\theta|)$ was
expressed in degrees. Within-cohort permutation tests for this alignment
($B = 2000$ permutations) were performed separately for each cohort.

\subsection{Genus-level linear models and FDR correction}
\label{sec:fdr-methods}

Three genus-level hypothesis families were tested, each on all 62 CLR
genera, using the 62-feature harmonised CLR matrix as input
(Table~\ref{tab:fdr-summary}).

\paragraph{Family 1 --- PD main effect.}
For each genus, a linear model was fitted:
\[
y_g \sim \mathrm{Cohort} + \mathrm{PD}
\]
where $y_g$ is the CLR abundance of genus $g$, Cohort is a categorical
factor with Finland as reference, and PD is a binary factor (0 = control,
1 = PD). The PD coefficient ($\hat\beta_{\mathrm{PD}}$, standard error, $t$
statistic, and $p$-value) was extracted from each model.

\paragraph{Family 2 --- PD-by-cohort interaction.}
For each genus, a full interaction model was fitted:
\[
y_g \sim \mathrm{Cohort} \times \mathrm{PD}
\]
A joint $F$-test comparing the interaction model against the additive model
(Family 1) was used to obtain a single $p$-value per genus representing
evidence that the PD effect differed across cohorts.
Cohort-specific PD effect estimates ($\hat\beta_{\mathrm{PD,Finland}}$,
$\hat\beta_{\mathrm{PD,Malaysia}}$, $\hat\beta_{\mathrm{PD,USA}}$) were
extracted from the full interaction model.

\paragraph{Family 3 --- Cross-cohort heterogeneity.}
For each genus, a per-cohort linear model was fitted:
\[
y_g \sim \mathrm{PD} \quad (\text{within each cohort separately})
\]
Per-cohort PD estimates were combined by fixed-effect inverse-variance
meta-analysis ($\hat\beta_{\mathrm{FE}}$, $\hat{se}_{\mathrm{FE}}$,
$z$-statistic, $p_{\mathrm{FE}}$). Between-cohort heterogeneity was
quantified by Cochran's $Q$ ($Q$, $p_Q$) and $I^2$.

\paragraph{BH correction.}
P-values from each family were adjusted independently using the
Benjamini--Hochberg method~\cite{Benjamini1995}. Significance was declared
at $q < 0.05$. Both nominal $p$-values and BH-adjusted $q$-values are
reported. PC1 alignment permutation $p$-values from the three within-cohort
tests were also BH-corrected within that family (three tests).

All analyses were implemented in base R (version 4.5.2) using the script
\texttt{scripts/apply\_fdr\_genus\_portability.R}.

\subsection{Machine-learning classifiers}

Three classifiers were trained on the 62-feature CLR matrix: elastic-net
logistic regression (\texttt{glmnet}), random forest (\texttt{ranger}), and
XGBoost. Elastic-net regularisation ($\lambda$) was selected by 10-fold
cross-validation on the training set. Random forest used 500 trees.
XGBoost used a maximum depth of 6 and 100 boosting rounds. All classifiers
were evaluated under two validation schemes.

\paragraph{Pooled cross-validation.}
10-fold stratified repeated cross-validation (10 repeats) on the full
683-sample dataset. Performance metrics: AUROC, AUPRC, balanced accuracy,
sensitivity, and specificity.

\paragraph{Leave-one-cohort-out validation.}
Each cohort was held out as the test set in turn; the remaining two cohorts
formed the training set. LOCO is strict: no test-cohort samples influenced
training or hyperparameter selection.

\subsection{Biomarker portability gap}

$\mathrm{BPG}_m = \hat{\mathrm{perf}}_{\mathrm{pooled},m} -
\hat{\mathrm{perf}}_{\mathrm{LOCO},m}$ for $m \in \{\mathrm{AUROC},
\mathrm{AUPRC}, \mathrm{balanced\ accuracy}\}$. A large BPG indicates that
pooled performance overestimates portable predictive ability.

\subsection{Domain-correction feasibility}

DEBIAS-M~\cite{Shi2022} was assessed for compatibility with strict LOCO
testing. Transductive cohort-centring was evaluated as a leakage-free
alternative: within the held-out test cohort, each CLR feature was
mean-centred prior to prediction. Performance was assessed on AUROC,
sensitivity, and specificity.

\subsection{Feature-stability analysis}

Feature importance was extracted from each LOCO training scenario (three
held-out cohorts $\times$ three models = nine scenarios). For each model,
pairwise comparison of top-10 and top-20 feature sets across the three LOCO
scenarios yielded: (i) mean pairwise overlap count (top 10); (ii) mean
pairwise Jaccard similarity (top 20).

\subsection{Reproducibility}

All analyses used fixed random seeds (\texttt{set.seed(1)} in R; equivalent
in Python). Code is available at [repository URL]. The complete pipeline can
be rerun from raw MOTHUR count outputs.

% ============================================================
\section{Results}
\label{sec:results}

\subsection{Dataset characteristics}

The harmonised dataset comprised 683 individuals across three cohorts:
Finland ($n = 148$; 50.0\% PD), Malaysia ($n = 198$; 51.0\% PD), and the
United States ($n = 337$; 60.5\% PD). Sixty-two genus-level CLR features
were retained after cross-cohort taxonomic harmonisation.

\subsection{Cohort structure dominates compositional space}
\label{sec:geometry}

PCA of the CLR feature matrix showed that PC1 alone explained 61.2\% of
total variance. Cohort means were well separated along PC1: Malaysia had
a mean score of $+14.95$ (SD 2.05), Finland $+0.15$ (SD 0.83), and USA
$-8.81$ (SD 1.04), confirming that between-cohort differences constitute
the primary source of compositional variation.

PERMANOVA of the Aitchison distance matrix attributed $R^2 = 0.7505$ to
cohort identity (75.05\% of total variance) and $R^2 = 0.0029$ to PD
status (0.29\%; Table~\ref{tab:permanova}). The CDDR was therefore
$0.7505 / 0.0029 = 258.7$: cohort identity explains approximately 259 times
as much compositional variance as disease status.

The disease-effect vector $\mathbf{d}$ (mean CLR profile of PD subjects
minus that of controls, pooled across cohorts) showed a cosine similarity
of 0.067 with PC1, corresponding to an angle of 86.1$^\circ$. This near-
orthogonality indicates that the direction of greatest between-group
microbiome variation and the direction of PD-associated microbiome change
are geometrically distinct. Within-cohort permutation tests for PC1
alignment yielded $p_{\mathrm{perm}} \geq 0.252$ in all three cohorts
after BH correction ($q \geq 0.756$; Table~\ref{tab:pc-align}), consistent
with the null hypothesis that the PD-effect vector is not aligned with PC1
in any individual cohort.

\begin{table}[h]
\centering
\caption{PERMANOVA results (marginal model; 999 permutations;
Aitchison distance).}
\label{tab:permanova}
\begin{tabular}{lccc}
\toprule
Factor  & $R^2$  & $p$ & CDDR \\
\midrule
Cohort  & 0.7505 & 0.001 & --- \\
PD      & 0.0029 & 0.002 & 258.7 \\
\bottomrule
\end{tabular}
\end{table}

\subsection{Genus-level PD effects are cohort-dependent and do not form
a portable signature}
\label{sec:genus}

We tested all 62 CLR genera for three hypothesis families
(Table~\ref{tab:fdr-summary}).

\paragraph{PD main effect (Family 1).}
Of 62 genera tested in the cohort-adjusted linear model
($y_g \sim \mathrm{Cohort} + \mathrm{PD}$), 24 showed a nominal PD
association ($p < 0.05$) and 20 survived BH correction ($q < 0.05$;
Table~\ref{tab:main-effect-top}). The strongest associations were for
an unclassified genus cluster that was lower in PD (Unclassified:
$\hat\beta_{\mathrm{PD}} = -0.495$, SE = 0.083, $q = 2.38 \times 10^{-7}$)
and for the Gammaproteobacteria unclassified group, which was higher in PD
($\hat\beta_{\mathrm{PD}} = +0.484$, SE = 0.086, $q = 7.73 \times 10^{-7}$).
Of the 20 genera surviving FDR, 12 showed higher CLR abundance in PD and
8 showed lower abundance (Table~\ref{tab:main-effect-top}).

\paragraph{PD-by-cohort interaction (Family 2).}
Of 62 genera tested by joint $F$-test of the interaction model versus the
additive model, 32 showed nominal evidence of interaction ($p < 0.05$) and
25 survived BH correction ($q < 0.05$; Table~\ref{tab:interaction-top}).
The strongest interaction was for the unclassified genus cluster
($q = 1.93 \times 10^{-18}$), which showed a positive PD association in
Finland ($\hat\beta_{\mathrm{PD,Finland}} = +0.114$) but a negative
association in Malaysia ($\hat\beta_{\mathrm{PD,Malaysia}} = -1.663$) and
was near null in the USA ($\hat\beta_{\mathrm{PD,USA}} = -0.056$). The
second strongest interaction was for \textit{Oscillibacter}
($q = 3.16 \times 10^{-17}$; $\hat\beta_{\mathrm{PD,Finland}} = +0.944$,
$\hat\beta_{\mathrm{PD,Malaysia}} = -0.014$,
$\hat\beta_{\mathrm{PD,USA}} = -0.025$). These results indicate that the
direction and magnitude of genus-level PD effects depend strongly on which
cohort is considered.

\paragraph{Cross-cohort heterogeneity (Family 3).}
Per-cohort linear models were combined by fixed-effect meta-analysis for
each genus. Of 62 genera, 12 showed a consistent direction of effect
surviving BH correction on the meta-analysed $z$-statistic ($q_{\mathrm{FE}}
< 0.05$; Table~\ref{tab:het-top}). All 12 were higher in PD in the pooled
meta-analytic direction. The most consistent effect was for
Clostridiales\_unclassified ($\hat\beta_{\mathrm{FE}} = +0.210$, SE = 0.032,
$q_{\mathrm{FE}} = 2.50 \times 10^{-9}$) and
Lachnospiraceae\_unclassified ($\hat\beta_{\mathrm{FE}} = +0.172$,
SE = 0.029, $q_{\mathrm{FE}} = 1.78 \times 10^{-7}$).

However, Cochran's $Q$ test for heterogeneity revealed substantial
between-cohort variability in PD effects: 31 of 62 genera showed significant
heterogeneity after BH correction ($q_Q < 0.05$). Among the 12 genera with
significant meta-analytic PD effects, most also showed high $I^2$ values,
including Coriobacteriaceae\_unclassified ($I^2 = 0.809$,
$q_Q = 5.30 \times 10^{-3}$), Firmicutes\_unclassified ($I^2 = 0.880$,
$q_Q = 2.42 \times 10^{-4}$), and Erysipelotrichaceae\_unclassified
($I^2 = 0.924$, $q_Q = 2.09 \times 10^{-6}$). The notable exception was
Lachnospiraceae\_unclassified ($I^2 = 0.000$, $p_Q = 0.377$), which
showed a small but consistent PD-associated increase across all three
cohorts.

Taken together, these results indicate that genus-level disease effects
exist, but are predominantly cohort-conditioned. The PD-by-cohort
interaction is significant for 25 of 62 genera and the between-cohort
heterogeneity is significant for 31 of 62 genera. Only Lachnospiraceae
unclassified showed a consistent cross-cohort PD signal without significant
heterogeneity.

\begin{table}[h]
\centering
\caption{FDR correction summary by hypothesis family. BH: Benjamini--Hochberg.
All 62 genera tested in each family. PC1 alignment tests are cohort-level
($n = 3$~cohorts), not genus-level.}
\label{tab:fdr-summary}
\begin{tabular}{lcccccc}
\toprule
Family & Model & $n$ tested & Nominal & BH $q<0.05$ &
Min $p$ & Min $q$ \\
\midrule
PD main effect & $y \sim \mathrm{Cohort}+\mathrm{PD}$ &
  62 & 24 & 20 & $3.84\times10^{-9}$ & $2.38\times10^{-7}$ \\
PD$\times$Cohort interaction & Joint $F$-test &
  62 & 32 & 25 & $3.12\times10^{-20}$ & $1.93\times10^{-18}$ \\
Heterogeneity FE ($p_{\mathrm{FE}}$) & FE meta $z$ &
  62 & 12 & 12 & $4.04\times10^{-11}$ & $2.50\times10^{-9}$ \\
Heterogeneity ($p_Q$) & Cochran $Q$ &
  62 & 35 & 31 & --- & --- \\
PC1 alignment & Permutation & 3 & 0 & 0 & 0.252 & 0.756 \\
\bottomrule
\end{tabular}
\end{table}

\begin{table}[h]
\centering
\caption{Top 10 genera by PD main effect (Family 1; $y_g \sim
\mathrm{Cohort} + \mathrm{PD}$). BH $q$-values calculated across all 62
genera. Direction: higher/lower CLR abundance in PD relative to controls.}
\label{tab:main-effect-top}
\begin{tabular}{lrrrrl}
\toprule
Genus & $\hat\beta_{\mathrm{PD}}$ & SE &
  $p$ & $q$ & Direction \\
\midrule
Unclassified                     & $-0.495$ & 0.083 & $3.84\times10^{-9}$ & $2.38\times10^{-7}$ & lower  \\
Gammaproteobacteria\_unclassified&  $+0.484$ & 0.086 & $2.50\times10^{-8}$ & $7.73\times10^{-7}$ & higher \\
Odoribacter                      &  $+0.462$ & 0.086 & $1.10\times10^{-7}$ & $2.27\times10^{-6}$ & higher \\
Megasphaera                      & $-0.579$ & 0.110 & $2.15\times10^{-7}$ & $3.33\times10^{-6}$ & lower  \\
Prevotella                       & $-0.273$ & 0.057 & $2.14\times10^{-6}$ & $2.65\times10^{-5}$ & lower  \\
Clostridiales\_unclassified      &  $+0.471$ & 0.100 & $2.83\times10^{-6}$ & $2.93\times10^{-5}$ & higher \\
Streptococcus                    &  $+0.365$ & 0.081 & $7.01\times10^{-6}$ & $6.21\times10^{-5}$ & higher \\
Oscillibacter                    &  $+0.193$ & 0.045 & $2.31\times10^{-5}$ & $1.79\times10^{-4}$ & higher \\
Paraprevotella                   & $-0.414$ & 0.100 & $3.84\times10^{-5}$ & $2.64\times10^{-4}$ & lower  \\
Erysipelotrichaceae\_unclassified&  $+0.307$ & 0.078 & $8.66\times10^{-5}$ & $5.37\times10^{-4}$ & higher \\
\bottomrule
\end{tabular}
\end{table}

\begin{table}[h]
\centering
\caption{Top 10 genera by PD-by-cohort interaction (Family 2; joint
$F$-test of $y_g \sim \mathrm{Cohort}\times\mathrm{PD}$ vs additive model).
Cohort-specific $\hat\beta_{\mathrm{PD}}$ from the interaction model;
Finland is the reference cohort. BH $q$-values across 62 genera.}
\label{tab:interaction-top}
\begin{tabular}{lrrrrrr}
\toprule
Genus & $p_{\mathrm{joint}}$ & $q_{\mathrm{joint}}$ &
  $\hat\beta_{\mathrm{FIN}}$ & $\hat\beta_{\mathrm{MYS}}$ &
  $\hat\beta_{\mathrm{USA}}$ \\
\midrule
Unclassified                     & $3.12\times10^{-20}$ & $1.93\times10^{-18}$ & $+0.114$ & $-1.663$ & $-0.056$ \\
Oscillibacter                    & $1.02\times10^{-18}$ & $3.16\times10^{-17}$ & $+0.944$ & $-0.014$ & $-0.025$ \\
Gammaproteobacteria\_unclassified& $9.03\times10^{-16}$ & $1.87\times10^{-14}$ & $+0.192$ & $+1.551$ & $-0.037$ \\
Prevotella                       & $3.64\times10^{-15}$ & $5.63\times10^{-14}$ & $-1.125$ & $+1.551$ & $-0.037$ \\
Acidaminococcus                  & $4.43\times10^{-8}$  & $5.49\times10^{-7}$  & $-0.803$ & --- & --- \\
Escherichia\_Shigella            & $3.92\times10^{-7}$  & $4.05\times10^{-6}$  & $-0.683$ & --- & --- \\
Coprococcus                      & $2.45\times10^{-6}$  & $2.17\times10^{-5}$  & $-0.893$ & --- & --- \\
Odoribacter                      & $4.65\times10^{-6}$  & $3.61\times10^{-5}$  & $+1.173$ & --- & --- \\
Lactobacillales\_unclassified    & $9.13\times10^{-5}$  & $6.29\times10^{-4}$  & $+0.192$ & --- & --- \\
Lactococcus                      & $1.14\times10^{-4}$  & $7.08\times10^{-4}$  & $-0.525$ & --- & --- \\
\bottomrule
\multicolumn{6}{l}{\footnotesize FIN = Finland; MYS = Malaysia; USA = USA.
Cohort betas not shown for ranks 5--10 to conserve space;} \\
\multicolumn{6}{l}{\footnotesize full values in \texttt{results/fdr/genus\_disease\_by\_cohort\_interactions\_fdr.csv}.} \\
\end{tabular}
\end{table}

\begin{table}[h]
\centering
\caption{Top 10 genera by fixed-effect meta-analytic PD association
(Family 3). $I^2$ quantifies between-cohort heterogeneity; $p_Q$ is
Cochran's $Q$ $p$-value. All 12 genera with $q_{\mathrm{FE}} < 0.05$ are
shown; all show higher CLR abundance in PD. BH $q$-values across 62 genera.}
\label{tab:het-top}
\begin{tabular}{lrrrrrr}
\toprule
Genus & $\hat\beta_{\mathrm{FE}}$ & SE & $q_{\mathrm{FE}}$ &
  $Q$ & $p_Q$ & $I^2$ \\
\midrule
Clostridiales\_unclassified      & $+0.210$ & 0.032 & $2.50\times10^{-9}$  & 8.12  & 0.017 & 0.754 \\
Lachnospiraceae\_unclassified    & $+0.172$ & 0.029 & $1.78\times10^{-7}$  & 1.95  & 0.377 & 0.000 \\
Coriobacteriaceae\_unclassified  & $+0.147$ & 0.029 & $9.10\times10^{-6}$  & 10.48 & 0.005 & 0.809 \\
Firmicutes\_unclassified         & $+0.156$ & 0.031 & $9.10\times10^{-6}$  & 16.66 & 0.0002& 0.880 \\
Bacteroidales\_unclassified      & $+0.155$ & 0.032 & $2.07\times10^{-5}$  & 19.86 & 0.0001& 0.899 \\
Erysipelotrichaceae\_unclassified& $+0.115$ & 0.026 & $6.46\times10^{-5}$  & 26.15 & $2.1\times10^{-6}$ & 0.924 \\
Bacteroidetes\_unclassified      & $+0.136$ & 0.031 & $1.05\times10^{-4}$  & 19.92 & $4.7\times10^{-5}$ & 0.900 \\
Fusobacteriaceae\_unclassified   & $+0.132$ & 0.031 & $1.24\times10^{-4}$  & 22.27 & $1.5\times10^{-5}$ & 0.910 \\
Gammaproteobacteria\_unclassified& $+0.108$ & 0.025 & $1.34\times10^{-4}$  & 48.46 & $3.0\times10^{-11}$& 0.959 \\
Bacteria\_unclassified           & $+0.125$ & 0.030 & $2.43\times10^{-4}$  & 23.61 & $7.5\times10^{-6}$ & 0.915 \\
\bottomrule
\end{tabular}
\end{table}

\subsection{Pooled validation overestimates microbiome biomarker portability}
\label{sec:ml}

Three classifiers were evaluated under pooled repeated cross-validation
and LOCO validation (Table~\ref{tab:performance}). Under pooled
cross-validation, all models showed apparently useful discrimination, with
AUROC values of 0.810 (elastic net), 0.881 (random forest), and 0.919
(XGBoost). AUPRC values were 0.821, 0.894, and 0.926, respectively.

Performance decreased markedly under LOCO validation
(Table~\ref{tab:performance}). Mean LOCO AUROC was 0.600 for elastic net,
0.515 for random forest, and 0.547 for XGBoost. AUROC portability gaps were
0.210, 0.366, and 0.371. AUPRC portability gaps were 0.173, 0.321, and
0.325. The models with the highest pooled AUROC showed the largest loss under
LOCO testing. Cohort-specific LOCO results revealed pronounced heterogeneity
(Table~\ref{tab:loco-by-cohort}): no model achieved AUROC above 0.76 across
all three held-out cohort scenarios simultaneously.

\begin{table}[h]
\centering
\caption{Pooled cross-validation and mean LOCO performance with biomarker
portability gaps (BPG). AUROC: area under the ROC curve; AUPRC: area under
the precision--recall curve.}
\label{tab:performance}
\begin{tabular}{lcccccc}
\toprule
Model & Pooled & LOCO & BPG & Pooled & LOCO & BPG \\
      & AUROC  & AUROC & (AUROC) & AUPRC & AUPRC & (AUPRC) \\
\midrule
Elastic net   & 0.810 & 0.600 & 0.210 & 0.821 & 0.648 & 0.173 \\
Random forest & 0.881 & 0.515 & 0.366 & 0.894 & 0.573 & 0.321 \\
XGBoost       & 0.919 & 0.547 & 0.371 & 0.926 & 0.601 & 0.325 \\
\bottomrule
\end{tabular}
\end{table}

\begin{table}[h]
\centering
\caption{LOCO AUROC by held-out cohort.}
\label{tab:loco-by-cohort}
\begin{tabular}{lccc}
\toprule
Model & Finland held out & Malaysia held out & USA held out \\
\midrule
Elastic net   & 0.418 & 0.757 & 0.623 \\
Random forest & 0.525 & 0.370 & 0.649 \\
XGBoost       & 0.562 & 0.416 & 0.664 \\
\bottomrule
\end{tabular}
\end{table}

\subsection{Domain-adaptive correction does not recover portable
classification}

DEBIAS-M was assessed for compatibility with strict LOCO evaluation.
Under the available workflow, DEBIAS-M required phenotype metadata from
the held-out cohort during reference-step fitting, constituting held-out
phenotype-label leakage and rendering it incompatible with a strict LOCO
design. DEBIAS-M was therefore not applied.

Transductive cohort-centring was applied as a leakage-free alternative:
within the held-out test cohort, each CLR feature was mean-centred prior to
prediction. This produced degenerate classifiers across all three models:
AUROC = 0.500, sensitivity = 1.000, specificity = 0.000
(Table~\ref{tab:corrected}). The classifiers predicted PD for every test
subject, indicating that removing cohort-centred structure eliminated the
apparent discriminative signal. This result is consistent with the
interpretation that cohort-structured compositional variation was the primary
driver of pooled predictive performance.

\begin{table}[h]
\centering
\caption{Performance of classifiers trained on transductively cohort-centred
features, under LOCO validation.}
\label{tab:corrected}
\begin{tabular}{lccc}
\toprule
Model & AUROC & Sensitivity & Specificity \\
\midrule
Elastic net   & 0.500 & 1.000 & 0.000 \\
Random forest & 0.500 & 1.000 & 0.000 \\
XGBoost       & 0.500 & 1.000 & 0.000 \\
\bottomrule
\end{tabular}
\end{table}

\subsection{Feature importance is unstable across cohort-held-out training
regimes}

Feature importance was extracted from each LOCO training scenario and
compared across held-out-cohort conditions (Table~\ref{tab:stability}).
Stability was low across all classifiers. Mean top-10 overlap was 2.3
features for elastic net, 3.0 for random forest, and 2.0 for XGBoost, all
well below the maximum possible overlap of 10. Mean top-20 Jaccard
similarity was 0.334, 0.380, and 0.225, respectively. Top-ranked features
varied substantially depending on which cohort was withheld from training,
reinforcing the interpretation that the apparent biomarker sets were
cohort-conditioned rather than portable.

\begin{table}[h]
\centering
\caption{Feature stability across LOCO training regimes (three held-out
cohort scenarios per model).}
\label{tab:stability}
\begin{tabular}{lccc}
\toprule
Model & Mean top-10 overlap & Mean top-20 overlap & Mean Jaccard (top 20) \\
\midrule
Elastic net   & 2.3 & 10.0 & 0.334 \\
Random forest & 3.0 & 11.0 & 0.380 \\
XGBoost       & 2.0 &  7.3 & 0.225 \\
\bottomrule
\end{tabular}
\end{table}

% ============================================================
\section{Discussion}
\label{sec:discussion}

This study presents a cross-cohort diagnostic workflow for detecting
microbiome biomarker non-portability under compositional cohort structure.
Applied to three PD gut microbiome cohorts, the workflow identified a
consistent failure mode: disease-associated genus-level effects were
statistically detectable but did not form a portable biomarker signature.
Cohort background dominated Aitchison-space variance (CDDR = 258.7),
the disease-effect vector was nearly orthogonal to the dominant
compositional axis (86.1$^\circ$), pooled machine-learning performance
overestimated portability, LOCO validation reduced AUROC toward the range
0.515--0.600, and feature importance was unstable across held-out-cohort
regimes.

The genus-level results clarify the mechanism. Twenty of 62 genera showed
a cohort-adjusted PD association surviving FDR correction, confirming that
disease-associated signals exist within this dataset. However, 25 of 62
genera showed significant PD-by-cohort interaction, and 31 of 62 showed
significant between-cohort heterogeneity in the PD effect direction or
magnitude. This degree of cohort-dependence means that no single genus-level
signature was consistent enough across cohorts to support portable
classification. The transductive cohort-centring result reinforces this: when
cohort-structured variation was removed, classifiers had no remaining
discriminative signal.

The pooled validation failure is therefore not merely a power issue. The
62-feature CLR matrix is rich enough to detect 20 FDR-significant PD
associations under cohort control. The problem is that these associations are
predominantly cohort-conditioned. Random cross-validation distributes
cohort-specific structure across training and test folds, producing optimistic
estimates. LOCO validation prevents this by requiring the disease signal to
generalise to a compositionally different cohort background.

The biomarker portability gap quantified this discrepancy directly. Models
with the highest pooled AUROC (XGBoost: 0.919) showed the largest portability
gap (0.371). Elastic net, the most interpretable model, showed a smaller gap
(0.210) and more stable feature selection, which may reflect its stronger
regularisation and consequently lower sensitivity to cohort-specific
structure.

The within-cohort heterogeneity of genus-level PD effects also has practical
implications for meta-analysis. The 12 genera with FDR-significant pooled FE
effects would appear in a meta-analysis as replicated disease associations.
However, most of these also show high $I^2$ values (range: 0.754--0.959 for
10 of the 12 genera), indicating that the apparent replication is driven by a
consistent direction of effect rather than consistent magnitude. The exception,
Lachnospiraceae\_unclassified ($I^2 = 0$, $q_{\mathrm{FE}} = 1.78 \times
10^{-7}$), represents the clearest candidate for a portable PD-associated
microbial signal in this dataset, but its small effect size
($\hat\beta_{\mathrm{FE}} = +0.172$ CLR units) and the absence of additional
replication in the LOCO prediction results suggest caution before nominating
it as a biomarker.

The proposed workflow is not specific to Parkinson's disease. The five
diagnostic steps --- (1) CDDR and variance partitioning, (2) disease-axis
geometry, (3) LOCO prediction, (4) BPG estimation, and (5) feature-stability
analysis --- can be applied to any multi-cohort microbiome biomarker dataset.
A CDDR near or above 50 is a prior indicator that pooled validation will
overestimate portable predictive ability. The disease-axis angle provides a
complementary geometric diagnostic. Together, these two pre-modelling checks
can flag datasets where LOCO portability testing is likely to reveal a large
portability gap.

Limitations include the restriction to three cohorts and 16S rRNA amplicon
sequencing. The 62-feature genus-level harmonisation used here is
conservative; species-level or strain-level analysis might reveal narrower
signals with different portability properties. The failure of DEBIAS-M in a
strict LOCO design may not generalise to all batch-correction methods; future
evaluation of reference-free or unsupervised correction approaches that do
not require held-out phenotype labels would be valuable. Finally, the PD
cohorts used here differ in geographic origin, which conflates microbial
biogeography with disease status to an unknown degree.

% ============================================================
\section{Conclusions}
\label{sec:conclusions}

Apparent microbiome biomarker performance can arise from cohort-structured
compositional variation rather than portable disease signal. In a three-cohort
Parkinson's disease dataset (683 individuals, 62 harmonised CLR features),
cohort structure explained 75.05\% of Aitchison-space variance while disease
status explained 0.29\%. Twenty of 62 genera showed FDR-significant
cohort-adjusted PD associations, but 25 of 62 showed significant
PD-by-cohort interaction and 31 of 62 showed significant between-cohort
heterogeneity. Pooled classification AUROC reached 0.919 but collapsed to a
mean LOCO AUROC of 0.547 for the same model (portability gap: 0.371).
Feature importance was unstable across cohort-held-out training regimes, and
transductive domain correction eliminated all discriminative signal. The
proposed five-step workflow --- compositional geometry, cohort-held-out
prediction, portability-gap estimation, feature-stability analysis, and
genus-level FDR-corrected heterogeneity testing --- provides a practical
diagnostic strategy to identify microbiome biomarker non-portability before
candidate signatures are interpreted as transportable across clinical cohorts.

% ============================================================
\section*{Declarations}

\textbf{Competing interests:} The authors declare no competing interests.

\textbf{Data availability}

All sequencing datasets analysed in this study are publicly available from
the European Nucleotide Archive (ENA) and the NCBI Sequence Read Archive
(SRA) under the following BioProject identifiers: PRJNA494620 (Malaysia),
PRJEB4927 (Finland), and PRJEB14674 (United States). Corresponding study
identifiers are SRP199959, ERP004264, and ERP016332, respectively.

No new sequencing data were generated. The harmonised genus-level abundance
matrix, processed data tables, and all R scripts required to reproduce the
analyses and figures are publicly available at Zenodo:
\url{https://doi.org/10.5281/zenodo.18644851}.

% ============================================================
\begin{thebibliography}{99}

\bibitem{Aitchison1982}
Aitchison J (1982) The statistical analysis of compositional data.
\textit{J R Stat Soc Ser B} 44:139--177.

\bibitem{Anderson2001}
Anderson MJ (2001) A new method for non-parametric multivariate analysis
of variance. \textit{Austral Ecol} 26:32--46.

\bibitem{Barichella2019}
Barichella M et al.\ (2019) Unraveling gut microbiota in Parkinson's
disease and atypical parkinsonism. \textit{Mov Disord} 34:396--405.

\bibitem{Bedarf2017}
Bedarf JR et al.\ (2017) Functional implications of microbial and viral
gut metagenome changes in early stage L-DOPA-naive Parkinson's disease
patients. \textit{Genome Med} 9:39.

\bibitem{Benjamini1995}
Benjamini Y, Hochberg Y (1995) Controlling the false discovery rate:
a practical and powerful approach to multiple testing.
\textit{J R Stat Soc Ser B} 57:289--300.

\bibitem{Oksanen2022}
Oksanen J et al.\ (2022) vegan: Community Ecology Package.
R package version 2.6-4.

\bibitem{Sampson2016}
Sampson TR et al.\ (2016) Gut microbiota regulate motor deficits and
neuroinflammation in a model of Parkinson's disease.
\textit{Cell} 167:1469--1480.

\bibitem{Shi2022}
Shi P et al.\ (2022) DEBIAS-M: Domain-adaptive microbiome batch correction.
\textit{bioRxiv} 2022.08.02.502474.

\end{thebibliography}

\end{document}
